\expandafter\gdef\csname SNside1\endcsname{\node[anchor=north west,inner sep=0pt,text width=48mm,font=\SideFont,align=left] at ([xshift=138mm,yshift=-618.99pt]current page.north west) {\textcolor{SNink}{\textbf{错项排除 · 实时性的判据是期限}}\par\vspace{1.5mm}控制任务必须在五毫秒内响应，即使某机器多数时候只用一毫秒，偶尔延迟十毫秒也可能违反硬实时要求。排除``实时系统就是运行速度最快的系统''时，应指出缺少最坏响应与截止期限这一保证。};}
\expandafter\gdef\csname SNside2\endcsname{\node[anchor=north west,inner sep=0pt,text width=48mm,font=\SideFont,align=left] at ([xshift=138mm,yshift=-95.00pt]current page.north west) {\textcolor{SNink}{\textbf{最小反例 · 单核也有并发}}\par\vspace{1.5mm}一个CPU先运行A十毫秒，再运行B十毫秒，二者在同一时间段交替推进，已满足并发，却没有在同一时刻执行两条用户指令。把``系统有多个任务''直接推出``任务必并行''，就混淆了时间区间与时间点。\par\smallskip\begin{tikzpicture}[x=1mm,y=1mm,>=Stealth,node distance=5mm]
\node[draw=SNink!50,fill=SNblue,rounded corners=1mm,text width=34mm,align=center,inner sep=1.2mm,font=\fontsize{7.5}{9}\selectfont] (m0) at (0,0) {A:10ms};
\node[draw=SNink!50,fill=SNblue,rounded corners=1mm,text width=34mm,align=center,inner sep=1.2mm,font=\fontsize{7.5}{9}\selectfont] (m1) at (0,-11) {B:10ms};
\node[draw=SNink!50,fill=SNblue,rounded corners=1mm,text width=34mm,align=center,inner sep=1.2mm,font=\fontsize{7.5}{9}\selectfont] (m2) at (0,-22) {单核交替};
\draw[->,SNink] (m0) -- (m1);
\draw[->,SNink] (m1) -- (m2);
\end{tikzpicture}
\par{\fontsize{7.5}{9}\selectfont 同一时间段交替推进是并发，不推出同一时刻指令并行。}};}
\expandafter\gdef\csname SNside3\endcsname{\node[anchor=north west,inner sep=0pt,text width=48mm,font=\SideFont,align=left] at ([xshift=138mm,yshift=-95.00pt]current page.north west) {\textcolor{SNink}{\textbf{最小反例 · 内核也要做普通加法}}\par\vspace{1.5mm}内核计算页框号或更新计数器时，会执行普通算术指令。它们在用户态通常也合法，所以``内核态执行的都是特权指令''不成立。限制方向是特权操作要求高权限，不能把这个蕴含关系反过来。};}
\expandafter\gdef\csname SNside4\endcsname{\node[anchor=north west,inner sep=0pt,text width=48mm,font=\SideFont,align=left] at ([xshift=138mm,yshift=-144.70pt]current page.north west) {\textcolor{SNink}{\textbf{课程连接 · 与组成原理共用一次入口}}\par\vspace{1.5mm}组成原理说明硬件如何保存断点、改变模式并定位入口；操作系统接着说明入口代码保存哪些通用寄存器以及怎样分派服务。做题先划出硬件响应与软件处理的边界，再判断某一步由谁完成。};}
\expandafter\gdef\csname SNside5\endcsname{\node[anchor=north west,inner sep=0pt,text width=48mm,font=\SideFont,align=left] at ([xshift=138mm,yshift=-95.00pt]current page.north west) {\textcolor{SNink}{\textbf{错项排除 · 系统调用不保证阻塞}}\par\vspace{1.5mm}读取的数据若已经在内核缓存且可立即返回，系统调用仍跨越了保护边界，却可能始终由原进程完成。判断阻塞要看所需事件是否尚未满足，判断进入内核要看是否通过受控入口，两项不能互相替代。};}
\expandafter\gdef\csname SNside6\endcsname{\node[anchor=north west,inner sep=0pt,text width=48mm,font=\SideFont,align=left] at ([xshift=138mm,yshift=-95.00pt]current page.north west) {\textcolor{SNink}{\textbf{不变量 · 先判范围，再求和}}\par\vspace{1.5mm}界限寄存器按长度定义时，合法逻辑地址始终满足\(0\le a<L\)。即使\(B+a\)仍落在已安装的物理内存内，也不能放行\(a=L\)。这个不变量保护的是当前进程的合法区间，与整台机器是否还有空间无关。};}
\expandafter\gdef\csname SNside7\endcsname{\node[anchor=north west,inner sep=0pt,text width=48mm,font=\SideFont,align=left] at ([xshift=138mm,yshift=-95.00pt]current page.north west) {\textcolor{SNink}{\textbf{错项排除 · 结构优势有适用条件}}\par\vspace{1.5mm}微内核减少高权限功能、分离服务，有利于故障隔离与扩展；一次服务却可能多出消息传递和保护域切换。比较表中的``有利于''是结构机制，不是对任意实际系统性能、安全性的无条件排名。};}
\expandafter\gdef\csname SNside8\endcsname{\node[anchor=north west,inner sep=0pt,text width=48mm,font=\SideFont,align=left] at ([xshift=138mm,yshift=-95.00pt]current page.north west) {\textcolor{SNink}{\textbf{课程连接 · 开机结构与磁盘格式不同}}\par\vspace{1.5mm}中断入口、页表和调度队列属于本次运行所需结构；分区表、根目录和磁盘索引节点通常是持久存储的一部分。把题中对象按``随运行建立''与``跨重启保留''定位，才能判断初始化是否需要重新创建。};}
\expandafter\gdef\csname SNside9\endcsname{\node[anchor=north west,inner sep=0pt,text width=48mm,font=\SideFont,align=left] at ([xshift=138mm,yshift=-95.00pt]current page.north west) {\textcolor{SNink}{\textbf{错项排除 · 客户内核不是主机最高权限}}\par\vspace{1.5mm}虚拟机里的客户操作系统可以认为自己处于内核态，但敏感操作仍受VMM与硬件虚拟化机制控制。客户能管理其虚拟资源，不代表它能任意改写主机物理资源；必须分清客户视图和实际保护边界。};}
\expandafter\gdef\csname SNside10\endcsname{\node[anchor=north west,inner sep=0pt,text width=48mm,font=\SideFont,align=left] at ([xshift=138mm,yshift=-95.00pt]current page.north west) {\textcolor{SNink}{\textbf{课程连接 · 指针自身与所指对象分开}}\par\vspace{1.5mm}数据结构中的链表头指针可能是函数局部变量，所指结点却由动态分配得到；函数退出与结点释放属于不同生存期问题。分析内存图时为``指针变量''和``动态对象''分别画位置，能避免把整个链表误判为栈对象。};}
\AddToShipoutPictureFG{\begin{tikzpicture}[remember picture,overlay]\ifcsname SNside\thepage\endcsname\csname SNside\thepage\endcsname\fi\end{tikzpicture}}